% Nouveau EGA V (cote 25): archive cote 25, site page 98 (PDF page 99).
% typescript. Draft compared: none.
% AI-transcribed by Claude Opus 5.5 (claude-opus-5-5, adaptive thinking, effort high), Message Batches API, batch ega25b (msgbatch_01A7n5YuXubUPNK1zsepi8RG), 2026-11-13; not yet checked by eye.
\scanpage{25/98}
\mspage{1}
\hand{Nouveau EGA V}
\var{\inserted{Faisceaux inversibles et diviseurs sur les fibrés projectifs et multiprojectifs. Systèmes linéaires de diviseurs. \var{\deleted{familles de}}{}{deletion}{ega25_25_098-v0002}}}{Faisceaux inversibles et diviseurs sur les fibrés projectifs et multiprojectifs. Systèmes linéaires de diviseurs. \var{\deleted{familles de}}{}{deletion}{ega25_25_098-v0002}}{insertion}{ega25_25_098-v0001}
% NOTE: a long handwritten diagonal stroke runs across the page from the struck title down to ``de L.'' (end of Corollaire 1.3); a second stroke crosses the line ``Une autre façon de formuler'' towards the lower left corner

IV 2\var{\over{0}{1}}{1}{writing over}{ega25_25_098-v0003} \var{\deleted{\uline{Compléments sur les faisceaux inversibles}, \uline{les} \uline{diviseurs de}}}{}{deletion}{ega25_25_098-v0004}
\var{\deleted{\var{\over{C}{\unsure{+}}}{\unsure{+}}{writing over}{ega25_25_098-v0006}\uline{artier}, \uline{les systêmes linéaires de diviseurs}.}}{}{deletion}{ega25_25_098-v0005}

1. \uline{Détermination des faisceaux inversibles sur les fibrés projectifs}.
\uline{Application\var{\xout{x}}{}{typewriter strike-out}{ega25_25_098-v0007}} \uline{au \var{\deleted{groupe}}{}{deletion}{ega25_25_098-v0008} des}\hand{/}\uline{automorphismes d'un fibré projectif}.

\uline{Théorème} 1.1. Soient S un préschéma, \underline{E} un Module loc libre sur \var{\xout{S qui}}{}{typewriter strike-out}{ega25_25_098-v0009} de
rang \var{\inserted{\unsure{fini}}}{\unsure{fini}}{insertion}{ega25_25_098-v0010} $\geqslant 2$ en tout point,
\var{\xout{n'est nul en aucun point,}}{}{typewriter strike-out}{ega25_25_098-v0011} P = P(\underline{E}) le fibré projectif qu'il définit. Alors
\var{\margin{*}}{\placedmargin{*}}{marginal note placed}{ega25_25_098-v0012} pour tout Module inversible \underline{L} sur P, on peut trouver une famille
\var{\xout{\illegible{1}}}{}{typewriter strike-out}{ega25_25_098-v0013} $(S_n)_{n\in\underline{Z}}$ d'ouverts disjoints de S recouvrant S, indexée par \underline{Z}, et
\var{\xout{\unsure{pour tout n}$\in$\underline{Z}}}{}{typewriter strike-out}{ega25_25_098-v0014} un Module inversible \var{\xout{$\underline{L}_n$}}{}{typewriter strike-out}{ega25_25_098-v0015}\var{\inserted[typed]{$\underline{M}$}}{$\underline{M}$}{insertion}{ega25_25_098-v0016} sur S\var{\xout{$_n$}}{}{typewriter strike-out}{ega25_25_098-v0017} , detelle façon que
\var{\deleted{\var{\xout{\unsure{$\underline{L}|p^{-1}(S_n)$ soit isomorphe à $p_n^{-1}(\underline{L}_n)\otimes\underline{O}$ la restriction de}}}{}{typewriter strike-out}{ega25_25_098-v0019}}}{}{deletion}{ega25_25_098-v0018} la restricti-%
on de \underline{L} à $f^{-1}(S_n)$ soit isomorphe à celle de $\underline{M}\otimes_{\underline{O}_S}\underline{O}_P(1) = f^*(\underline{M})(n)$.
De plus, la \var{\xout{famille}}{}{typewriter strike-out}{ega25_25_098-v0020} famille des $S_n$ \var{\xout{\unsure{et la classe}}}{}{typewriter strike-out}{ega25_25_098-v0021} est uniquement déter-%
minée par ses conditions, ainsi que \underline{M} à isomorphisme unique près.

\uline{Remarque} 1.2. Lorsque on ne fait pas d'hypothèse sur le rang de \underline{E}, S se
décompose canoniquement en \var{\inserted{le}}{le}{insertion}{ega25_25_098-v0022}\var{\inserted[typed]{préschéma}}{préschéma}{insertion}{ega25_25_098-v0023}somme de trois ouverts $S^0,S^1,S^2$ tel que sur
ceux-ci, le rang de \underline{E} soit resp. 0, 1, et $\geqslant 2$. Alors la détermination des
Modules inversibles sur P se ramène à celle des $f^{-1}(S^{i})$ pour i=0,1,2.
Le cas i=2 est justiciable de 1.1. , d'autre part $f^{-1}(S^1)$ est $S^1$-isomor-%
phe à $S^1$ donc son groupe de Picard n'est autre que Pic($S^1$), enfin
$f^{-1}(S^0)$ est vide donc son groupe de Picard est nul.

\uline{Corollaire} 1.3. Sous les conditions de 1.1. supposons S connexe\var{\inserted[typed]{non vide}}{non vide}{insertion}{ega25_25_098-v0024}. Alors
tout Module inversible \underline{L} sur P est isomorphe à un Module de la forme
$f^*(\underline{M})(n)$, où $n\in\underline{Z}$ et \underline{M} est un Module inversible sur S. De plus n est
uniquement déterminé, et \underline{M} est déterminé à isom unique près, par la donnée
de \underline{L}.

Une autre façon de formuler ce corollaire est la suivante. Considé-%
rons les homomorphismes naturels Pic(S)$\longrightarrow$Pic(\unsure{P}), \underline{Z} $\longrightarrow$Pic(P), le premier
\var{\margin{98}}{\placedmargin{98}}{marginal note placed}{ega25_25_098-v0025}
